% CP-VI-0030
% target_chapter: VI/20 — Gluing Affine Schemes
% source_unit_id: IMP-DL-0B4CC237FA-SPR000660-87506BC7
% source: Downloads/theory-of-algebraic-geometry-3.tex L660-L719
% solution_provenance: SOURCE_EXPLICIT_SOLUTION

\subsection*{Problem 14. Gluing schemes}

Let \(\{X_i\}_{i\in I}\) be a family of schemes. Suppose that for each \(i\neq j\), we are given an open subscheme
\[
U_{ij}\subseteq X_i
\]
and an isomorphism of schemes
\[
\varphi_{ij}:U_{ij}\to U_{ji}.
\]
Assume that the following compatibility conditions hold:

\begin{enumerate}[label=\textup{(\arabic*)}]
	\item For each \(i\neq j\),
	\[
	\varphi_{ji}=\varphi_{ij}^{-1}.
	\]
	
	\item For all \(i,j,k\),
	\[
	\varphi_{ij}(U_{ij}\cap U_{ik})
	=
	U_{ji}\cap U_{jk},
	\]
	and
	\[
	\varphi_{ik}
	=
	\varphi_{jk}\circ \varphi_{ij}
	\]
	on \(U_{ij}\cap U_{ik}\).
\end{enumerate}

Show that there exists a scheme \(X\), together with morphisms
\[
\psi_i:X_i\to X,
\]
such that:

\begin{enumerate}[label=\textup{(\arabic*)}]
	\item each \(\psi_i\) is an isomorphism of \(X_i\) with an open subscheme of \(X\);
	
	\item the open subsets \(\psi_i(X_i)\) cover \(X\);
	
	\item
	\[
	\psi_i(U_{ij})
	=
	\psi_i(X_i)\cap \psi_j(X_j);
	\]
	
	\item
	\[
	\psi_i
	=
	\psi_j\circ \varphi_{ij}
	\]
	on \(U_{ij}\).
\end{enumerate}

\textbf{Solution.}

The idea is to construct \(X\) by taking the disjoint union of the schemes \(X_i\) and identifying points that correspond under the transition isomorphisms \(\varphi_{ij}\). Then we glue the structure sheaves using the same transition isomorphisms.

For convenience, we set
\[
U_{ii}=X_i,
\qquad
\varphi_{ii}=\operatorname{id}_{X_i}.
\]
This allows us to write the compatibility conditions uniformly.

First construct the underlying set.

Let
\[
Y=\coprod_{i\in I}X_i
\]
be the disjoint union of the underlying sets of the \(X_i\). We define an equivalence relation \(\sim\) on \(Y\) as follows.

For \(x\in X_i\) and \(y\in X_j\), declare
\[
x\sim y
\]
if
\[
x\in U_{ij}
\]
and
\[
y=\varphi_{ij}(x).
\]
The compatibility assumptions guarantee that this relation is symmetric and transitive.

Indeed, symmetry follows from
\[
\varphi_{ji}=\varphi_{ij}^{-1}.
\]
For transitivity, suppose
\[
x\in X_i,
\qquad
y=\varphi_{ij}(x)\in X_j,
\qquad
z=\varphi_{jk}(y)\in X_k.
\]
Then the cocycle condition gives
\[
z
=
\varphi_{jk}(\varphi_{ij}(x))
=
\varphi_{ik}(x),
\]
whenever the expression is defined. Hence
\[
x\sim z.
\]
Thus the identifications are consistent.

Define
\[
X=Y/{\sim}.
\]
Let
\[
\psi_i:X_i\to X
\]
be the natural map sending a point to its equivalence class.

We now put a topology on \(X\). A subset
\[
V\subseteq X
\]
is declared to be open if and only if
\[
\psi_i^{-1}(V)\subseteq X_i
\]
is open for every \(i\in I\). This is the quotient topology with respect to the maps \(\psi_i\).

With this topology, each \(\psi_i\) is continuous by construction.

We now show that \(\psi_i\) identifies \(X_i\) homeomorphically with an open subset of \(X\).

First, the image
\[
\psi_i(X_i)\subseteq X
\]
is open. Indeed, for every \(j\), we compute
\[
\psi_j^{-1}\bigl(\psi_i(X_i)\bigr)=U_{ji}.
\]
This is open in \(X_j\). Hence \(\psi_i(X_i)\) is open in \(X\).

Second, the map
\[
\psi_i:X_i\to \psi_i(X_i)
\]
is a homeomorphism. Its inverse is well-defined because the gluing identifications do not identify two distinct points inside the same \(X_i\). This follows from the cocycle condition with \(\varphi_{ii}=\operatorname{id}_{X_i}\).

Furthermore, the topology on \(\psi_i(X_i)\) agrees with the topology of \(X_i\), because for any open subset \(W\subseteq X_i\), the image \(\psi_i(W)\) is open in \(\psi_i(X_i)\), and conversely every open subset of \(\psi_i(X_i)\) pulls back to an open subset of \(X_i\).

The images \(\psi_i(X_i)\) cover \(X\), since every equivalence class in \(X\) contains a point from some \(X_i\). Thus
\[
X=\bigcup_{i\in I}\psi_i(X_i).
\]

Next we describe the overlaps.

By construction, a point of \(X_i\) and a point of \(X_j\) have the same image in \(X\) exactly when they are identified by \(\varphi_{ij}\). Hence
\[
\psi_i(X_i)\cap \psi_j(X_j)
=
\psi_i(U_{ij})
=
\psi_j(U_{ji}).
\]
Therefore
\[
\boxed{
	\psi_i(U_{ij})
	=
	\psi_i(X_i)\cap \psi_j(X_j).
}
\]
Also, for every \(x\in U_{ij}\),
\[
\psi_i(x)
=
\psi_j(\varphi_{ij}(x)).
\]
Thus
\[
\boxed{
	\psi_i=\psi_j\circ\varphi_{ij}
	\quad\text{on }U_{ij}.
}
\]

It remains to construct the structure sheaf on \(X\).

For an open subset \(V\subseteq X\), define
\[
V_i=\psi_i^{-1}(V)\subseteq X_i.
\]
A section of \(\mathcal O_X\) over \(V\) will be a family of sections
\[
(s_i)_{i\in I},
\qquad
s_i\in \Gamma(V_i,\mathcal O_{X_i}),
\]
which are compatible on overlaps.

More precisely, define
\[
\mathcal O_X(V)
\]
to be the set of all families
\[
(s_i)_{i\in I}\in \prod_{i\in I}\Gamma(V_i,\mathcal O_{X_i})
\]
such that for every pair \(i,j\),
\[
s_i|_{V_i\cap U_{ij}}
=
\varphi_{ij}^{\#}
\left(
s_j|_{V_j\cap U_{ji}}
\right).
\]
Here
\[
\varphi_{ij}^{\#}:
\mathcal O_{X_j}|_{U_{ji}}
\longrightarrow
(\varphi_{ij})_*\mathcal O_{X_i}|_{U_{ij}}
\]
is the isomorphism of sheaves induced by the isomorphism of schemes
\[
\varphi_{ij}:U_{ij}\to U_{ji}.
\]

Restriction maps for \(\mathcal O_X\) are defined componentwise:
if
\[
W\subseteq V,
\]
then
\[
(s_i)_{i\in I}
\longmapsto
(s_i|_{W_i})_{i\in I},
\]
where
\[
W_i=\psi_i^{-1}(W).
\]
The compatibility condition is preserved under restriction, so this gives a presheaf of rings.

This presheaf is in fact a sheaf. Indeed, the sheaf condition can be checked componentwise on the open cover
\[
X=\bigcup_i \psi_i(X_i),
\]
and each \(\mathcal O_{X_i}\) is already a sheaf. The cocycle condition for the maps \(\varphi_{ij}\) guarantees that sections which agree pairwise on overlaps glue uniquely.

Therefore we obtain a sheaf of rings
\[
\mathcal O_X
\]
on \(X\).

Now we show that each \(X_i\) is identified with an open subscheme of \(X\).

Restrict \(\mathcal O_X\) to the open subset \(\psi_i(X_i)\). Since
\[
\psi_i:X_i\to \psi_i(X_i)
\]
is a homeomorphism, we may pull back the restricted sheaf to \(X_i\). By construction of \(\mathcal O_X\), this pullback is canonically isomorphic to \(\mathcal O_{X_i}\). Therefore
\[
\psi_i:
(X_i,\mathcal O_{X_i})
\longrightarrow
(\psi_i(X_i),\mathcal O_X|_{\psi_i(X_i)})
\]
is an isomorphism of ringed spaces.

Moreover, the stalks of \(\mathcal O_X\) are local rings. Indeed, if \(x\in X\), then \(x\in \psi_i(X_i)\) for some \(i\). Choose \(x_i\in X_i\) with
\[
\psi_i(x_i)=x.
\]
Since \(\psi_i\) identifies \(X_i\) with the open subset \(\psi_i(X_i)\), we have
\[
\mathcal O_{X,x}\cong \mathcal O_{X_i,x_i}.
\]
But \(X_i\) is a scheme, so \(\mathcal O_{X_i,x_i}\) is a local ring. Hence \(\mathcal O_{X,x}\) is local.

Thus
\[
(X,\mathcal O_X)
\]
is a locally ringed space.

Finally, \(X\) is a scheme because it is covered by the open subsets
\[
\psi_i(X_i),
\]
and each open subset \(\psi_i(X_i)\) is isomorphic to the scheme \(X_i\). Since being a scheme is local on the underlying topological space, \(X\) is a scheme.

We have constructed a scheme \(X\) and morphisms
\[
\psi_i:X_i\to X
\]
such that:

\begin{enumerate}[label=\textup{(\arabic*)}]
	\item each \(\psi_i\) identifies \(X_i\) with an open subscheme of \(X\);
	
	\item the open subschemes \(\psi_i(X_i)\) cover \(X\);
	
	\item
	\[
	\psi_i(U_{ij})
	=
	\psi_i(X_i)\cap \psi_j(X_j);
	\]
	
	\item
	\[
	\psi_i
	=
	\psi_j\circ \varphi_{ij}
	\quad\text{on }U_{ij}.
	\]
\end{enumerate}

Therefore
\[
\boxed{
	\text{schemes can be glued along compatible open subscheme isomorphisms.}
}
\]

The glued scheme \(X\) is unique up to unique isomorphism with these properties.
